\abstractpagehu{Kivonat}{%
  Jelen diplomadolgozat az automatikus hangszerfelismerés problémáját vizsgálja
  zenei felvételek alapján, modern gépi tanulási módszerek alkalmazásával.
  A kidolgozott rendszer valós időben képes osztályozni hét hangszerkategóriát
  (gitár, zongora, ének, vonós, nádas fúvós, rézfúvós és környezeti zaj)
  monofón audiojel alapján.

  A megvalósított megoldás három spektrális jellemző típust (Log-Mel
  spektrogram, STFT és MFCC) használ, amelyek kétdimenziós konvolúciós neurális
  hálózatba (2D~CNN) kerülnek betáplálásra. A modell tanítása Google Colab
  környezetben történt. Az adathalmazt saját tiszta internetes minták és
  mikrofonnal/telefonnal újrafelvételezett minták, valamint az ESC-50
  környezeti zajadatbázis egyesítésével hoztam létre. Az adataugmentációs modul
  a valós felvételi körülményeket szimulálja visszhang, környezeti zaj és
  mikrofonkarakterisztika hozzáadásával.

  A kísérleti eredmények versenyképes osztályozási pontosságot mutatnak a
  tesztadatokon, és a rendszert valós idejű mikrofonos felismeréssel is
  validáltam. A dolgozat hozzájárul annak demonstrálásához, hogy elérhető
  számítási erőforrásokkal is megvalósítható egy robusztus hangszerosztályozó
  rendszer.
}{%
\textit{\textbf{Kulcsszavak:}} hangszerfelismerés, konvolúciós neurális hálózat, Mel-spektrogram, STFT, MFCC, audioosztályozás, hullámforma, spektrális kép%
}

